% ======================================================================
% 2.3 Сферические координаты. Ответственный: D — Филиппов Андрей.
%
% Что сдать (main.tex: условие 555–591, порядок 592–620, пример 708–750;
% требования к отчёту 1862–1870):
%   1. Изображение полушара: радиус, плоскость z=0.
%   2. Обоснование пределов r, theta, phi; матрица частных производных
%      x,y,z по r,theta,phi; определитель; проверка |J|=r^2*sin(phi).
%   3. Повторный интеграл от g(r,theta,phi) с явными пределами; аналитическое решение.
%   4. Метод средних прямоугольников при n=8,16,32,64;
%      шаги h_r=R/n, h_theta=2*pi/n, h_phi=pi/(2n).
%   5. Проверка объёма: интегрировать r^2*sin(phi), сравнить с V=2*pi*R^3/3.
%   6. Таблица значений и ошибок, график ошибки от h=1/n, вывод о сходимости.
%
% Параметры варианта 3 (приложение А): R=1.
% Код: solution.py (блок D), запуск python experiments.py 2.3.
% Рисунки: figures/2-3-*.pdf.
% ======================================================================

\subsection{Сферические координаты}

\taskhead{Исходные данные и аналитическое решение}
Для варианта 3 задан $R=1$. Рассматривается интеграл
\[
I_{\mathrm{с}}=\iiint_{\Omega_{\mathrm{с}}}z^2\,dx\,dy\,dz,
\qquad
\Omega_{\mathrm{с}}=\{x^2+y^2+z^2\leq R^2,\ z\geq0\}.
\]
\todo{D: описать область; обосновать пределы $r,\theta,\varphi$; выписать замену координат; составить матрицу частных производных и вычислить определитель (проверить $|J|=r^2\sin\varphi$); записать функцию $g(r,\theta,\varphi)$ и повторный интеграл; вычислить аналитически (пример: main.tex, строки 708–750).}

\taskhead{Программный результат и сравнение}
\todo{D: листинг вызова своей функции из solution.py; указать шаги $h_r=R/n$, $h_\theta=2\pi/n$, $h_\varphi=\pi/(2n)$.}
Аналитическое значение $I_{\mathrm{с}}$: \todo{D: подставить}.
\begin{table}[htbp]
\centering\small
\caption{Сравнение численного интеграла с аналитическим значением $I_{\mathrm{с}}$}
\label{tab:2-3}
\begin{tabular}{|r|r|r|r|r|}\hline
$n$ & $h$ & $I_n$ & $\Delta_n$ & $\delta_n$\\\hline
8 & — & — & — & —\\
16 & — & — & — & —\\
32 & — & — & — & —\\
64 & — & — & — & —\\
\hline\end{tabular}
\end{table}

\taskhead{Проверка объёма}
\todo{D: заменить исходную функцию на 1 (в новых координатах интегрировать $r^2\sin\varphi$); вычислить на всех сетках и сравнить с $V=2\pi R^3/3$; пояснить результат.}

\taskhead{Графическая проверка}
\taskfigure[0.6]{figures/2-3-region.pdf}{Верхняя половина шара при $R=1$ (равный масштаб осей)}{fig:2-3-region}
\taskfigure[0.75]{figures/2-3-error.pdf}{Абсолютная ошибка от нормированного шага $h=1/n$ (логарифмические оси)}{fig:2-3-error}
\todo{D: пояснить график.}

\taskhead{Вывод}
\todo{D: вывод о сходимости; объяснить поведение ошибки при удвоении $n$.}
